\section{Experimental setup}
\label{sec:experiments}
\label{sec:empirical_evaluation}

Evaluations on a single testbed have limited external validity~\cite{apruzzese2022modeling, apruzzese2023real, he2023adversarial, vitorino2023sok}, and detectors trained on one environment transfer poorly to another~\cite{dang2022detecting, dang2023learning, dang2026stfedxids}. We therefore evaluate on three corpora calibrated to the published marginal distributions of CSE-CIC-IDS2018 (enterprise)~\cite{sharafaldin2018toward}, UNSW-NB15 (contemporary hybrid)~\cite{moustafa2015unsw}, and CIDDS-001 (cloud OpenStack)~\cite{ring2017flow}. Each corpus is synthetic. It is generated to match its testbed's statistics and is not a replay of the original packet captures.

Each corpus contains 12,000 flows, of which 9,000 are benign and 3,000 malicious. The malicious flows form five families of 1,000, 500, 700, 500 and 300 flows: a SYN flood, a UDP flood, brute-force authentication against SSH and FTP, SQL injection and cross-site scripting, and low-and-slow command-and-control infiltration. The data are split into 8,000 training, 1,000 validation and 3,000 test flows. The threat model is an ingress perimeter that inspects forward telemetry before bidirectional session termination (Section~\ref{sec:formulation}).

The targets are XGBoost (CUDA histogram method) and LightGBM, each a 60-tree gradient-boosted ensemble of depth 6 with learning rate 0.1. Experiments ran on an NVIDIA RTX 5000 Ada GPU with 16 GB of memory, using Python 3.13.11, XGBoost 3.4.1, LightGBM 4.7.0, CatBoost 1.2.10 and scikit-learn 1.8.0.

We compare five evasion methods against FFE:
\begin{itemize}
    \item Naive coordinate descent in $\mathbb{R}^d$ with a fixed step $\Delta = 0.05\,\sigma_f$ and no immutability, monotonicity or rate-coupling checks.
    \item ZOO-NIDS~\cite{chen2017zoo}: zeroth-order optimization with randomly sampled coordinates, symmetric difference quotients, Adam updates and box bounds.
    \item HopSkipJump-NIDS~\cite{chen2020hopskipjump}: a decision-based boundary attack initialized from a benign flow.
    \item CAA/MOEVA~\cite{simonetto2024constrained, ghamizi2020moeva}: R-NSGA-III evolutionary search with population 50 and 25 generations.
    \item Traffic Manipulator PSO~\cite{han2021traffic}: particle swarm optimization over packet intervals and forward dummy bytes, with swarm size 50 and 37 iterations (up to 1,850 queries).
    \item FFE: the Lie-steering procedure (Algorithms~\ref{alg:ffe_master}--\ref{alg:tsbs}) with step budget $M = 20$.
\end{itemize}

ASR is the fraction of attacks with $P(\text{Malicious} \mid x^*) < 0.5$. PFR is the fraction of candidate flows that pass the protocol sanity checks, the non-negativity checks and the exact coupled rate equations, that is, $x^* \in \mathcal{M}_{\text{flow}}$. For unconstrained baselines that never produce an evasive candidate, PFR is computed over evasive candidates only. Real ASR counts an attack only when it both evades and is physically feasible:
\begin{equation}
    \text{Real ASR} = \frac{1}{N} \sum_{i=1}^N \mathbb{I}\left[ F(x_i^*) < 0 \;\land\; x_i^* \in \mathcal{M}_{\text{flow}} \right],
\end{equation}
where $F$ is the detector score and negative values denote evasion. Average queries counts the model evaluations per sample, and total time is the wall-clock time for the cohort. The perturbation cost is the mean feature-normalized $L_1$ distance over the $d=12$ features:
\begin{equation}
    \bar{D}_{L_1}(x^*, x) = \frac{1}{d} \sum_{j=1}^d \frac{|x_j^* - x_j|}{\sigma_j},
\end{equation}
where $\sigma_j$ is the empirical standard deviation of feature $j$ in the training corpus.

Each method is run on the same cohort of 100 correctly classified malicious flows sampled at random from the test partition. The per-family and cross-dataset cohorts also contain 100 flows each. The ablation and defense studies use balanced cohorts of 20 flows per attack family, so their figures differ from the master cohort (for example, 11.6 against 12.9 queries at $M=20$). Table~\ref{tab:transferability} also trains CatBoost and Random Forest models on the same corpus. Each study trains its own model instances, so clean accuracy for the same architecture differs slightly across tables (99.70\% in Table~\ref{tab:multidataset_benchmark}, 99.72\% in Table~\ref{tab:defense_evaluation}). All experiments use a single random seed, so we report no variance.
